
\section{Estimating the Curvature Factors}
\label{app:matnormal-w1}
\Cref{alg:ours} turns the single-batch update of
\Cref{sec:sample-loss-lmo} into an optimizer by averaging its two inputs,
the gradient in the linear objective and the second moment in the
outergradient constraint. Averaging over recent steps draws on more samples
while downweighting gradients computed farther from the current iterate. The factor-gradient
momenta $\widehat M_A,\widehat M_B$ supply
the projections used in the direction step~\eqref{eq:factor-directions},
while an EMA accumulates the vectors $p,q$ that define the diagonal
curvature factors $P,Q$. \Cref{sec:weighted-estimation-w1} derives these
quantities from LoRA factor gradients alone, and
\Cref{sec:scale-placement-w1} explains why the accumulated vectors are
normalized after averaging. The update is invariant to the scale of the
curvature factors, but the average that estimates them depends on that scale.

\subsection{Estimation from the Factor Gradients}
\label{sec:weighted-estimation-w1}

\paragraph{Running averages.} Let $s$ index optimizer steps, let $G_s$ be
the batch gradient at step $s$, and write $g_s=\vec(G_s)$.
\Cref{alg:ours} uses look-ahead momentum for the gradient input, matching
line~\ref{ln:mom},
\[
  M_{t+1}=\beta_1M_t+(1-\beta_1)G_t,
  \qquad
  \widehat M_t=\beta_1M_{t+1}+(1-\beta_1)G_t.
\]
Unrolling this recurrence gives objective weights $m_{t,s}$ on past
gradients,
\[
  m_{t,s}=\beta_1^2(1-\beta_1)\beta_1^{t-1-s}
  \quad(s<t),
  \qquad
  m_{t,t}=1-\beta_1^2 .
\]
For the curvature input it uses an exponential moving average, which keeps
memory constant and downweights older gradients computed farther from the
current iterate. We write the induced weights on past steps as
\[
  \EMA_{s \leq t}[\phi(G_s)]=\sum_{s\le t}w_{t,s}\,\phi(G_s),
  \qquad
  w_{t,s}=(1-\beta_2)\beta_2^{t-s}.
\]
Thus the curvature factors are fit to the average
\[
  \overline{\Sigma}_t=\EMA_{s\leq t}[g_sg_s^\top],
\]
rather than to the current batch alone. This subsection targets the
observed-gradient average; it does not recompute old gradients at the current
weights.

\paragraph{Using the factor gradients.} LoRA backpropagation gives
$G_{A,s}=B_s^\top G_s$ and
$G_{B,s}=G_sA_s^\top$~\eqref{e-grad-wrt-factors}. Forming the dense batch
gradient $G_s$ would require a $d_{\mathrm{out}}\times d_{\mathrm{in}}$
matrix per layer, so the optimizer builds the averages from the factor
gradients. The curvature estimator uses the following modeling steps for
the recent observed gradient stream and the changing LoRA factors.
\begin{enumerate}[leftmargin=1.6em,itemsep=2pt,topsep=2pt]
  \item Let
  $g_{s,i}=\vec(\nabla\ell_{x_{s,i}}(W_s))$ be the per-sample gradients in
  the size-$n_s$ minibatch at step $s$. Replace their second moment by the
  outer product of the batch gradient,
  \[
    \frac1{n_s}\sum_{i=1}^{n_s}g_{s,i}g_{s,i}^\top
    \;\leadsto\;
    g_sg_s^\top ,
  \]
  the batch-gradient second-moment approximation used by
  AdaGrad~\cite{adagrad}, Adam~\cite{adam}, and Shampoo~\cite{shampoo}.
  It is exact at batch size one, and the accumulated second moment is then
  $\overline{\Sigma}_t$.
  \item When relating the dense observed moment to factor-gradient moments,
  replacing $B_s^\top G_s$ by $B_t^\top G_s$ and
  $G_sA_s^\top$ by $G_sA_t^\top$ assumes weighted factor-projection
  stability. The concrete stability bounds used by the online fit are stated in
  \Cref{prop:online-stationary-w1}; it measures the whole quadratic-form
  change, including the cross terms created by changing the factors.
  \item Model the resulting local moment by diagonal Kronecker factors. We write
  $P^\star,Q^\star$ for these model factors to distinguish them from the
  estimates $P,Q$ in \Cref{alg:ours}. Since
  $Q^\star\otimes P^\star$ is unchanged by the
  rescaling $(P^\star,Q^\star)\mapsto(aP^\star,a^{-1}Q^\star)$, we fix
  $\norm{\diag P^\star}_\infty=\norm{\diag Q^\star}_\infty=1$ and let the
  scalar $\kappa_t>0$ carry the magnitude. The exact model is
  \begin{equation}\label{e-diag-kron-second-mom-model-w1}
    \EMA_{s \leq t}[g_sg_s^\top]
    =
    \kappa_t(Q^\star_t\otimes P^\star_t).
  \end{equation}
\end{enumerate}

\paragraph{Factor momenta.} The objective decomposes, under the linearized
product update~\eqref{eq:DeltaW-lin} with $\widehat M_t$ replacing the batch
gradient in~\eqref{eq:polorafull}, as
\[
  \ip{\widehat M_t,\,\Delta B\,A+B\,\Delta A}
  =\ip{\widehat M_tA^\top,\Delta B}+\ip{B^\top\widehat M_t,\Delta A},
\]
so the update depends on $\widehat M_t$ only through $\widehat M_tA^\top$
and $B^\top\widehat M_t$. The momenta $\widehat M_A,\widehat M_B$ in
\Cref{alg:ours} apply the same look-ahead rule to $B_s^\top G_s$ and
$G_sA_s^\top$. Replacing $A_s,B_s$ by $A_t,B_t$ in these first-moment sums
requires the corresponding $m$-weighted projection-replacement bounds
\[
\begin{aligned}
  \left\|
    \sum_{s\le t}m_{t,s}(B_s^\top G_s-B_t^\top G_s)
  \right\|_2
  &=
  O(\epsilon)
  \left\|
    \sum_{s\le t}m_{t,s}B_t^\top G_s
  \right\|_2,\\
  \left\|
    \sum_{s\le t}m_{t,s}(G_sA_s^\top-G_sA_t^\top)
  \right\|_2
  &=
  O(\epsilon)
  \left\|
    \sum_{s\le t}m_{t,s}G_sA_t^\top
  \right\|_2 .
\end{aligned}
\]

\paragraph{Moments of the factor gradients.} To estimate $P^\star$ and
$Q^\star$ from factor gradients, we link their second moments to $\overline{\Sigma}_t$.
Since
$\vec(G_{A,s})=(I\otimes B_s^\top)\vec(G_s)$,
\begin{align*}
 \Sigma_{A,t} &:=  \EMA_{s \leq t}\big[ \vec(G_{A,s})  \vec(G_{A,s})^\top\big]  \\
 &=  \EMA_{s \leq t}\big[ (I\otimes B_s^\top)\vec(G_s)\vec(G_s)^\top   (I\otimes B_s)\big]  \\
   &\approx  (I\otimes B_t^\top) \EMA_{s \leq t}\big[ g_s g_s^\top   \big]  (I\otimes B_t),
\end{align*}
where the last line uses $B_t$ in place of $B_s$ inside the weighted
average. Applying the diagonal Kronecker model~\eqref{e-diag-kron-second-mom-model-w1}
then gives
\begin{align*}
 \Sigma_{A,t}
   &\approx
   \kappa_t\,
   (I\otimes B_t^\top)
   (Q^\star_t\otimes P^\star_t)
   (I\otimes B_t) \\
   &= \kappa_t\,Q^\star_t\otimes(B_t^\top P^\star_tB_t),
\end{align*}
where the equality is the mixed-product rule of Kronecker products. The
analogous calculation for $B$ gives
\begin{align}
  \Sigma_{A,t}
    &\approx \kappa_t\,Q^\star_t\otimes C^\star_{A,t},
  \nonumber \\
  \Sigma_{B,t}
    &:=
      \EMA_{s \leq t}\big[
        \vec(G_{B,s}^\top)\vec(G_{B,s}^\top)^\top
      \big]
    \approx \kappa_t\,P^\star_t\otimes C^\star_{B,t},
  \label{eq:observable-moments-w1}
\end{align}
where $C^\star_{A,t}=B_t^\top P^\star_tB_t$ and
$C^\star_{B,t}=A_tQ^\star_tA_t^\top$ are the $r\times r$ matrices
of~\eqref{eq:ours-closure} evaluated at the model factors.

\paragraph{Fitting the factors.} We fit $P^\star$ and $Q^\star$ to the
moments in~\eqref{eq:observable-moments-w1}. In this paragraph all
quantities are at the current step, so we suppress the subscript $t$.
Following KL-Shampoo~\cite{klshampoo}, we use the log-determinant
divergence~\cite{logdetdiv}
\[
  D(\Sigma,S)=\Tr(\Sigma S^{-1})-\log\det(\Sigma S^{-1})-d,
\]
defined for positive definite matrices $\Sigma$ and $S$ of common dimension
$d$. Under the block model below, the $A$-gradient moment would recover
$Q^\star$ once a small $r\times r$ closure matrix is fixed, and the
$B$-gradient moment would recover $P^\star$ once its small closure matrix is
fixed. The fit therefore freezes the opposite-side closure,
\begin{equation}\label{eq:factor-fit-w1}
  \hat{q}(\widehat C_A)
  =\argmin_{q>0}\;
  D\big(\Sigma_A,\; \diag(q) \otimes \widehat C_A\big),
  \qquad
  \hat{p}(\widehat C_B)
  =\argmin_{p>0}\;
  D\big(\Sigma_B,\;\diag(p) \otimes \widehat C_B\big),
\end{equation}
with $\widehat C_A,\widehat C_B\succ0$. The blockwise first-order conditions
give the closed forms
\[
  \hat{q}(\widehat C_A)=\frac1r\diag\!\left(\EMA_{s \leq t}[G_{A,s}^\top \widehat C_A^{-1}G_{A,s}]\right),
  \qquad
  \hat{p}(\widehat C_B)=\frac1r\diag\!\left(\EMA_{s \leq t}[G_{B,s} \widehat C_B^{-1}G_{B,s}^\top]\right).
\]
For fixed positive definite $\Sigma$ and $S$, let
\[
  c^\star(\Sigma,S) :=\argmin_{c>0}D(\Sigma,cS)=\Tr(\Sigma S^{-1})/d
\]
be the best scalar multiple of $S$.

\begin{lemma}[Closure invariance of the fixed fit]\label{lem:consistency-w1}
Suppose the factor-gradient moments~\eqref{eq:observable-moments-w1} hold with
equality, with the model factors normalized so that
$\norm{\diag P^\star}_\infty=\norm{\diag Q^\star}_\infty=1$, and with all
diagonal entries of $P^\star,Q^\star$ positive. Assume also that the modeled
closures $C^\star_A$ and $C^\star_B$ are positive definite. For any positive
definite closures $\widehat C_A$ and $\widehat C_B$, the
fits~\eqref{eq:factor-fit-w1} return
\[
  \hat{q}(\widehat C_A)=\kappa\,c^\star\big(C^\star_A,\widehat C_A\big)\,\diag(Q^\star),
  \qquad
  \hat{p}(\widehat C_B)=\kappa\,c^\star\big(C^\star_B,\widehat C_B\big)\,\diag(P^\star),
\]
so $\hat{q}/\norm{\hat{q}}_\infty=\diag(Q^\star)$ and
$\hat{p}/\norm{\hat{p}}_\infty=\diag(P^\star)$.
\end{lemma}
\begin{proof}
Under~\eqref{eq:observable-moments-w1}, $\Sigma_A$ is
block diagonal with $j$th block $\kappa\,q^\star_j\,C^\star_A$, and the
candidate $\diag(q)\otimes\widehat C_A$ of~\eqref{eq:factor-fit-w1} at a
point $q>0$ is block diagonal with $j$th block $q_j\,\widehat C_A$. Traces
and log-determinants add across blocks, so the objective is
$\sum_jD(\kappa q^\star_jC^\star_A,\,q_j\widehat C_A)$, one term per block,
and the $j$th term is minimized at its best scalar multiple,
\[
\hat{q}_j=c^\star(\kappa q^\star_jC^\star_A,\widehat C_A)
=\kappa\,q^\star_j\,c^\star(C^\star_A,\widehat C_A).
\]
The factor $\kappa\,c^\star(C^\star_A,\widehat C_A)$ is the same for every
$j$, and $\norm{\diag Q^\star}_\infty=1$, so the normalization removes it.
The claim for $\hat{p}$ follows with the roles of the two factors exchanged.
\end{proof}

\paragraph{Online estimator.} The fit~\eqref{eq:factor-fit-w1} averages past
factor gradients with weights $w_{t,s}$. \methodname{} implements this
average online. Given raw EMA vectors $p_t,q_t$, the implementation forms
the effective diagonal metrics
\[
  D_P(p_t)=\diag\!\bigl(p_t/\norm{p_t}_\infty+\delta\mathbf 1\bigr),
  \qquad
  D_Q(q_t)=\diag\!\bigl(q_t/\norm{q_t}_\infty+\delta\mathbf 1\bigr),
\]
where $\delta>0$ is the relative damping parameter. The raw vectors
$p_t,q_t$ are initialized with positive entries, so these normalizations are
defined before the first curvature update. We state the inverse formulas for
nondegenerate steps where the damped closures below are positive definite:
\[
\begin{aligned}
  C_{A,t}&=B_t^\top D_P(p_t)B_t,
  &
  C_{B,t}&=A_tD_Q(q_t)A_t^\top,\\
  \widetilde C_{A,t}^{-1}
  &=(C_{A,t}+\delta\lambda_{\max}(C_{A,t})I)^{-1},
  &
  \widetilde C_{B,t}^{-1}
  &=(C_{B,t}+\delta\lambda_{\max}(C_{B,t})I)^{-1}.
\end{aligned}
\]
The Newton--Schulz routine computes the inverse square roots whose squares
give these inverse closures, up to numerical approximation. At step $t$ the
optimizer computes one damped fitted observation with the closures available
at that step, then folds that observation into the online recurrence
\begin{equation}\label{eq:ours-klcoupling-w1}
\begin{aligned}
  \hat{q}_t &= \frac1r\diag\!\left(G_{A,t}^\top \widetilde C_{A,t}^{-1}G_{A,t}\right),
  &
  \hat{p}_t &= \frac1r\diag\!\left(G_{B,t} \widetilde C_{B,t}^{-1}G_{B,t}^\top\right),
  \\
  q_{t+1} &= \beta_2q_t+(1-\beta_2)\hat{q}_t,
  &
  p_{t+1} &= \beta_2p_t+(1-\beta_2)\hat{p}_t,
  \\
  Q_{t+1} &= \diag\!\bigl(q_{t+1}/\norm{q_{t+1}}_\infty\bigr),
  &
  P_{t+1} &= \diag\!\bigl(p_{t+1}/\norm{p_{t+1}}_\infty\bigr).
\end{aligned}
\end{equation}
Unrolling the EMA gives
\[
  q_{t+1}
  =
  \beta_2^{t+1}q_0+q^\circ_{t+1},
  \qquad
  p_{t+1}
  =
  \beta_2^{t+1}p_0+p^\circ_{t+1},
  \qquad
  q^\circ_{t+1}=\sum_{s\le t}w_{t,s}\hat q_s,
  \qquad
  p^\circ_{t+1}=\sum_{s\le t}w_{t,s}\hat p_s .
\]
The proposition below compares these tail-free averages with the fitted
vectors obtained by freezing the factor maps and inverse closures at the
current model factors. The raw vectors $q_{t+1},p_{t+1}$ differ from
$q^\circ_{t+1},p^\circ_{t+1}$ by the displayed initialization terms.

\begin{proposition}[EMA tracking from weighted factor and closure stability]
\label{prop:online-stationary-w1}
At step $t$, let
$q^\star_t=\diag(Q^\star_t)$ and
$p^\star_t=\diag(P^\star_t)$ be positive vectors with
$\norm{q^\star_t}_\infty=\norm{p^\star_t}_\infty=1$, defined by the exact
dense moment model
\[
  \sum_{s\le t}w_{t,s}\,g_sg_s^\top
  =
  \kappa_t(Q^\star_t\otimes P^\star_t).
\]
Set
$C^\star_{A,t}=B_t^\top P^\star_tB_t$ and
$C^\star_{B,t}=A_tQ^\star_tA_t^\top$, and assume these closures are
positive definite. Define
\[
  \overline G_{A,s}=B_t^\top G_s,\qquad
  \overline G_{B,s}=G_sA_t^\top,\qquad
  H^A_t=(C^\star_{A,t})^{-1},\qquad
  H^B_t=(C^\star_{B,t})^{-1}.
\]
Let $K^A_s=\widetilde C_{A,s}^{-1}$ and
$K^B_s=\widetilde C_{B,s}^{-1}$. Assume the observed factor maps and the
online inverse closures are stable in the weighted quadratic forms
\[
\begin{aligned}
  \frac{1}{r\kappa_t}
  \left\|
    \sum_{s\le t}w_{t,s}
    \bigl(
      G_{A,s}^\top H^A_tG_{A,s}
      -\overline G_{A,s}^\top H^A_t\overline G_{A,s}
    \bigr)
  \right\|_2
  &=
  O(\epsilon),
  \\
  \frac{1}{r\kappa_t}
  \left\|
    \sum_{s\le t}w_{t,s}
    G_{A,s}^\top(K^A_s-H^A_t)G_{A,s}
  \right\|_2
  &=
  O(\epsilon),
  \\
  \frac{1}{r\kappa_t}
  \left\|
    \sum_{s\le t}w_{t,s}
    \bigl(
      G_{B,s}H^B_tG_{B,s}^\top
      -\overline G_{B,s}H^B_t\overline G_{B,s}^\top
    \bigr)
  \right\|_2
  &=
  O(\epsilon),
  \\
  \frac{1}{r\kappa_t}
  \left\|
    \sum_{s\le t}w_{t,s}
    G_{B,s}(K^B_s-H^B_t)G_{B,s}^\top
  \right\|_2
  &=
  O(\epsilon).
\end{aligned}
\]
Then, for $\epsilon$ below a fixed numerical constant,
\[
\begin{aligned}
  \left\|
    \frac{q^\circ_{t+1}}{\norm{q^\circ_{t+1}}_\infty}
    -
    q^\star_t
  \right\|_\infty
  &=
  O(\epsilon),
  \\
  \left\|
    \frac{p^\circ_{t+1}}{\norm{p^\circ_{t+1}}_\infty}
    -
    p^\star_t
  \right\|_\infty
  &=
  O(\epsilon).
\end{aligned}
\]
\end{proposition}
\begin{proof}
Define the ideal frozen fits
\[
\begin{aligned}
  q^{\mathrm{id}}_t
  &=
  \frac1r\diag\!\sum_{s\le t}w_{t,s}
  \overline G_{A,s}^\top H^A_t\overline G_{A,s},
  \\
  p^{\mathrm{id}}_t
  &=
  \frac1r\diag\!\sum_{s\le t}w_{t,s}
  \overline G_{B,s}H^B_t\overline G_{B,s}^\top .
\end{aligned}
\]
The exact dense moment model and the mixed-product rule give
$q^{\mathrm{id}}_t=\kappa_tq^\star_t$ and
$p^{\mathrm{id}}_t=\kappa_tp^\star_t$. Adding and subtracting the terms with
$H^A_t$ and $H^B_t$ in the online averages, the four stability bounds and
$\norm{\diag M}_\infty\le\norm{M}_2$ imply
\[
  \norm{q^\circ_{t+1}-q^{\mathrm{id}}_t}_\infty
  +
  \norm{p^\circ_{t+1}-p^{\mathrm{id}}_t}_\infty
  =
  O(\epsilon)\kappa_t .
\]
For any nonzero vector $u$ and perturbation $e$ satisfying
$\norm{e}_\infty\le\rho\norm{u}_\infty$ with $\rho<1$,
\[
  \left\|
    \frac{u+e}{\norm{u+e}_\infty}
    -
    \frac{u}{\norm{u}_\infty}
  \right\|_\infty
  \le
  \frac{2\rho}{1-\rho}.
\]
Applying this inequality with $u=\kappa_tq^\star_t$ and
$u=\kappa_tp^\star_t$ proves the claim.
\end{proof}

The stability bounds are the only assumptions specific to the online
recurrence. They say that replacing the current factor maps by the maps used
when each gradient arrived, and replacing the ideal inverse closures by the
online damped inverse closures, changes the weighted quadratic fits only by a
small relative amount.

\paragraph{Comparison with KL-Shampoo.} KL-Shampoo~\cite{klshampoo} also
fits Kronecker factors with the log-determinant divergence, and its moving
average has the same shape as~\eqref{eq:ours-klcoupling-w1}. The fit here
differs in three ways.
\begin{itemize}[leftmargin=1.5em,itemsep=1pt,topsep=2pt]
  \item \emph{No asymptotic distribution limit.} KL-Shampoo fits the second
  moment of a stationary gradient distribution. Here $\EMA_{s \leq t}$ is a
  finite weighted sum over realized gradients. The online statement is a
  deterministic perturbation statement for that weighted sum, not an
  infinite-sample limit.
  \item \emph{The fitted moments come from $G_A$ and $G_B$.} KL-Shampoo
  fits Kronecker factors to the second moment of the weight gradient $G$.
  Here $G$ is never formed. Each diagonal factor is fit from the second
  moment of its corresponding LoRA factor gradient, and the two fits are
  coupled through $C_A$ and $C_B$.
  \item \emph{Fixed-moment statement.} KL-Shampoo
  justifies its moving average as a stochastic proximal-gradient step on
  the divergence objective, correct in the limit. Here
  \Cref{lem:consistency-w1} is an algebraic statement for fixed moments.
  Under the deterministic moment model, using an inexact frozen
  closure only rescales the fitted vector, which the normalization removes.
  The online EMA is covered by the weighted factor-and-closure stability
  statement in \Cref{prop:online-stationary-w1}.
\end{itemize}

\paragraph{Recovering the update of \Cref{alg:ours}.} Substituting the
fitted model $\kappa_t(Q_t\otimes P_t)$ for the second moment and the
look-ahead momentum $\widehat M_t$ for the batch gradient gives the
preconditioned spectral LMO~\eqref{eq:polorafull}
(\Cref{lem:pseudogradient-spectral}). The direction and magnitude steps of
\Cref{sec:sample-loss-lmo}, with the factor momenta standing in for the two
dense momentum projections, give the update~\eqref{eq:polora-update}. In
\Cref{alg:ours}, the curvature factors update after the direction step, so
step $t$ uses the estimate from step $t-1$. This avoids forming and
inverting $C_A$ and $C_B$ twice per optimizer step.

\subsection{Normalization Placement}
\label{sec:scale-placement-w1}
\methodname{} stores raw accumulated vectors $q_{t+1},p_{t+1}$ and normalizes
them only when forming the effective metrics $D_Q(q_{t+1})$ and
$D_P(p_{t+1})$. It does not normalize each fitted observation
$\hat q_t,\hat p_t$ before adding it to the EMA.

\paragraph{Raw scale does not change the implemented closures.}
For any $a,b>0$,
\[
  D_Q(aq)=D_Q(q),
  \qquad
  D_P(bp)=D_P(p).
\]
The rank-$r$ closures $C_A,C_B$ and the relative-damped inverse closures
$\widetilde C_A^{-1},\widetilde C_B^{-1}$ are therefore unchanged by
positive rescaling of the raw EMA vectors. This is the implementation-level
scale invariance. The magnitude $\kappa_t$ in the diagonal Kronecker
model~\eqref{e-diag-kron-second-mom-model-w1} is removed before the
direction step sees the curvature factors.

\paragraph{Normalizing each observation would change the estimator.}
The recurrence in~\eqref{eq:ours-klcoupling-w1} estimates the weighted
arithmetic mean of the fitted observations,
\[
  q^\circ_{t+1}=\sum_{s\le t}w_{t,s}\hat q_s,
  \qquad
  p^\circ_{t+1}=\sum_{s\le t}w_{t,s}\hat p_s,
\]
with the raw vectors differing by the initialization terms in the unrolled
recurrence. Normalizing before the EMA would replace this by
\[
  \sum_{s\le t}w_{t,s}\frac{\hat q_s}{\norm{\hat q_s}_\infty},
  \qquad
  \sum_{s\le t}w_{t,s}\frac{\hat p_s}{\norm{\hat p_s}_\infty},
\]
which gives a small-magnitude and a large-magnitude fitted observation the
same maximum coordinate before averaging. The implemented order averages the
raw observations first and applies the projective normalization only when the
next effective metric is formed. The same average-then-normalize order appears
in normalized SGD~\cite{cutkosky_nigt}, where momentum precedes
normalization, and in Muon~\cite{denoise_ortho}, where momentum precedes
orthogonalization.
